
\section{Knowledge Graph \texorpdfstring{\emojikgtitle}{Knowledge Graph}}
\label{sec-kg}
A knowledge graph is a structured database that connects entities through well-defined relationships. It can either encompass a broad spectrum of general knowledge, such as the widely recognized Google Knowledge Graph~\cite{chen2020review, liu2020k}, or delve deeply into specialized domains, like the BioASQ dataset~\cite{tsatsaronis2015overview} for biomedical reasoning. The diverse information contained in a knowledge graph -- represented as entities, relationships, paths, and subgraphs -- serves as a valuable resource for enhancing various downstream tasks across different sectors, including question-answering~\cite{tian2024graph, wang2024knowledge, yasunaga2022deep}, commonsense reasoning~\cite{ilievski2021cskg}, fact-checking~\cite{kim2023factkg}, recommender systems~\cite{guo2020survey}, drug discovery~\cite{bonner2022review}, healthcare~\cite{chandak2023building}, and fraud detection~\cite{mao2022financial}.




\subsection{Application Tasks}
This section reviews representative applications that GraphRAG on KGs is used for.

\begin{itemize}[leftmargin=*]
    \item \textbf{Question-answering}: Question-answering (QA) can focus on a single domain or span across global knowledge. Typically, a query in text format is given, such as "What is the best way to predict a baby's eye color?" or "Were there fossil fuels in the ground when humans evolved?"~\cite{tian2024graph} -- the answer can be a sentence generated by a large language model (LLM), a selected text span from relevant documents, or even a specific choice in a multiple-choice QA scenario. In all these contexts, GraphRAG leverages knowledge graphs to retrieve relevant information, providing the necessary context or supporting facts to generate accurate answers.
    
    \item \textbf{Fact-Checking}: Fact-checking is to verify the truthfulness of statements by cross-referencing them with reliable sources of information. GraphRAG enhances this task by querying a knowledge graph to retrieve relevant facts and relational structures that either support or refute the given claim. GraphRAG identifies discrepancies or confirmations within the data by mapping the statement onto the knowledge graph, providing a thorough and evidence-based validation process.

    \item \textbf{Knowledge Graph Completion}: Knowledge graph completion is the task of predicting new facts to enhance the comprehensiveness of the graph and infer missing facts~\cite{zhang2023making}. GraphRAG addresses this task by retrieving structural knowledge around the triplets for inference, supplying essential structural knowledge, and enhancing the LLM inference.

    \item \textbf{Cybersecurity Analysis and Defense}: Cybersecurity Analysis and Defense aims to analyze and respond to vulnerabilities, weaknesses, attack patterns, and threat tactics. With the increasing complexity and volume of cybersecurity data, GraphRAG has been proposed to provide cybersecurity analysis with more comprehensive insights into potential attack vectors and mitigation strategies~\cite{rahman2024retrieval}.

\end{itemize}


\subsection{Knowledge Graph Construction}

We discuss how KGs are typically constructed. For each type of construction technique, we give examples of common KG databases. How a KG is constructed is important, as it can affect both its usefulness and function in different downstream tasks. We describe the main techniques below:
\begin{itemize}[leftmargin=*]
    \item {\bf Manual construction}: Some KGs are constructed manually via human annotation. WikiData~\cite{vrandevcic2014wikidata} is a KG that uses crowd-sourced efforts to gather a variety of knowledge. Each entity corresponds to a page in the Wikipedia encyclopedia. Another KG, the Unified Medical Language System (UMLS)~\cite{umls}, contains biomedical facts collected from numerous sources.

    \item {\bf Rule-based construction}: Many traditional approaches use rule-based techniques for constructing the graph. This takes the form of custom parsers and manually defined rules used to extract facts from raw text. Note that these parsers can differ depending on the source of the text. Prominent examples include ConceptNet~\cite{speer2017conceptnet}, which links different words together via assertions, and Freebase~\cite{bollacker2008freebase} which contains a wide variety of general facts. REANO~\cite{fang2024reano} extracts the entities and relations from a set of passages using traditional entity recognition (ER) and relation extraction (RE) methods, respectively. This includes SpaCy~\cite{honnibal2017spacy} for ER and TAGME~\cite{ferragina2010tagme} for RE. To extract the facts that connect two entities via a relation, they use DocuNet~\cite{zhang2021document}.

    \item {\bf LLM-based construction}: Recently, work has explored how LLMs can be used to construct KGs from a set of documents. In such a way, the LLM can automatically extract the entities and relations and link those together to form facts in the given text. Of note is that no ground-truth KG exists for these methods. Rather, they simply use a KG as a way to organize and represent a set of documents. For example, CuriousLLM~\cite{yang2024curiousllm} considers passages in the text as entities and determines whether two entities should be connected based on their encoded textual similarity. On the other hand, \citet{cheng2024structure} uses a manually-defined prompt to convert a piece of text into a KG. Graph-RAG~\cite{edge2024local} first divides each document into chunks and then uses an LLM to detect all the entities in each chunk, including their name, type, and description. To identify the relation between any two entities, both entities and a description of their relationship are passed to an LLM. An LLM is then used again to summarize the content of each entity and relation to arrive at their final title. Lastly, AutoKG~\cite{chen2023autokg} uses a combination of LLM embeddings and clustering techniques to construct a KG from a set of texts.
\end{itemize}




\subsection{Retriever}

Real-world facts in KGs can provide grounded information for generative models, enhancing the reliability of the model output. Given the structured nature of KGs, they are naturally well-suited for retrieval. The goal is for a given question or query to retrieve either relevant facts~\footnote{Throughout this paper, we will refer to facts as triples or edges, interchangeably} or entities that can help answer that question. Multiple considerations need to be considered during retrieval, including the type of facts we want to retrieve, the efficiency, and the amount of facts retrieved. In general, retrieval of KGs has two stages: identifying seed entities and retrieving facts or entities. We describe both below.


{\bf Identifying seed entities}: The first step in retrieving the relevant facts for a given query is to identify a set of ``seed entities'', which we'll refer to as $V_{\text{seed}}$. Seed entities are the initial entities that are chosen to be highly relevant to the original query. Given such, we expect that triples that contain any of these entities or are nearby in the graph should provide helpful context. Multiple techniques exist for identifying the seed entities. Some works~\cite{jiang2023structgpt, kim2023kg, gnn_llm, pullnet, zhang2022subgraph} assume that we are given a set of initial entities for each query. However, most works~\cite{feng2020scalable, sunthink, wen2023mindmap, zhanggreaselm, niu2024mitigating, yasunaga2022deep} attempt to extract the entities from the query. One approach is through entity extraction~\cite{al2020named}, which uses methods specifically designed for extracting entities from a given text. Most works only extract entities from the original query. Another common approach is to extract a set of entities that are semantically similar to the original query~\cite{wen2023mindmap, sanmartin2024kg}. HyKGE~\cite{jiang2024hykge} first generates a hypothesis and extracts entities from the original query and the hypothesis. Similarly, in order to reduce the possibility of hallucination, \citet{knowledgenavigator} uses an LLM to generate two similar questions and retrieves all entities found in the original and generated questions. In a similar vein, RoK~\cite{wang2024reasoning} first uses chain of thoughts reasoning to expand the original query, extracting the seed entities from the expanded query. 


{\bf Retrieval Methods}: The outcome of the previous step provides us with a set of entities that are related in some capacity to the query. These entities are then leveraged to retrieve a set of facts or entities that can aid us in answering the query. We summarize the core retrieval methods below.
\begin{itemize}[leftmargin=*]
    \item {\bf Traversal-based retriever}: These methods traverse the graph and extract paths to aid in answering a specific question. Given the set of seed entities, $V_{\text{seed}}$, \citet{yasunaga2021qa, yasunaga2022deep, zhanggreaselm} extract all paths up to length two between the entities in $V_{\text{seed}}$, resulting in a final entity set $V$. They further augment $V$ by including all triples that connect any two entities in $V$. Given $V$, both \cite{yasunaga2021qa, zhanggreaselm} only keep the top $k$ entities by relevance scores. This is calculated by training a separate model that takes the text embedding of the query and entity as input and outputs how relevant the entity is to the query. For \citet{yasunaga2022deep}, if $\lvert V \rvert > 200$, they randomly sample 200 entities. \citet{sunthink} use a version of beam-search to explore the KG. \citet{jiang2024hykge, feng2020scalable} extract all paths of length $k \leq 2$ between seed entities. Alternatively, LARK~\cite{choudhary2023complex} retrieves all facts that lie on paths $\leq k$ in length starting from the seed entities.     \citet{delile2024graph} first extract the shortest paths connecting all seed entities. They further prioritize some entities over others by considering the recency, importance, and relevance to the query of their associated text. OREOLM~\cite{hu2022empowering} traverse $k$ hops from the seed entities, contextualizing the importance of each relation and entity to a path via a learnable $d$-dimensional embedding and it's LM-encoded representation. \citet{zhang2022subgraph} introduce a trainable retriever that traverses the graph starting from each seed entity. They also train a model to score each newly visited edge, only keeping a portion of them. KG-RAG~\cite{sanmartin2024kg} works in a similar manner, scoring each edge by its relevance and similarity to the query via a dense retriever. They then use an LLM to decide which paths to explore in the next step.
    RoG~\cite{luoreasoning} uses instruction tuning to fine-tune an LLM to generate useful relation paths, which can be retrieved from the KG. 
    KnowledgeNavigator~\cite{knowledgenavigator} first uses the query to predict the expected number of hops, $h_Q$ needed to traverse in the retrieval stage. It then traverses $h_Q$ hops starting from the seed entities, using an LLM to score and prune irrelevant nodes. \citet{retrieve_rewrite_answer} operate in a similar manner; however, they choose which paths to traverse based solely on the relations. Furthermore, when scoring a path, all relations that lie on that path are considered when computing the score. RoK~\cite{wang2024reasoning} considers a different approach, using the Personalized PageRank (PPR) score to identify useful paths. They further augment these paths by including the 1-hop neighbors of the seed entities. PullNet~\cite{pullnet} assumes that each entity has an associated set of documents. Given a single seed entity, PullNet, traverses k hops, where in each iteration it extracts the facts for the newly observed entities. It also extracts any entities that are contained in documents associated with an entity found in the traversal. Furthermore, for each entity, only the top N facts are used, which are ranked via similarity to the query. KG-R3~\cite{pahuja2023retrieve} uses MINERVA~\cite{minerva}, a reinforcement learning approach to mining paths between entities, to retrieve a set of important paths between both entities in the fact. 
    \citet{wang2024knowledge} use an LLM to traverse the graph starting from the seed entities. At each iteration in the traversal, we choose the next node to visit by prompting an LLM. Specifically, given the information already collected in the traversal, the LLM is prompted to generate the remaining information needed to correctly answer the question. The neighboring node that best matches the required information is chosen as the next node to visit. They further instruction-tune the LLM. 


    \item {\bf Subgraph-based retriever}: These methods extract a subgraph of size $k$ around each of the seed entities. Facts that contain one of the seed entities or are nearby, should be highly relevant to answering the question. Furthermore, they may actually contain the answer itself. Each of \cite{niu2024mitigating, tian2024graph, jiang2023structgpt, kim2024causal} extract either the one or two hop subgraph around each seed entity. The final set of facts is the union of each individual subgraph. \citet{gao2022graph} propose to first extract the subgraph containing the seed entity and potential answers using the method in \citet{sun2018open}. This is then partitioned into a set of smaller subgraphs. Then, they design a framework to rank the subgraphs, keeping the top $k$ subgraphs for generation. For a question-choice pair, MVP-Tuning~\cite{mvp_tuning} considers the triple that contains the highest number of seed and choice entities. They further augment this by extracting the top $k$ most similar questions in the dataset using BM25~\cite{bm25}, and extract the triples for each of them.

    \item {\bf Rule-based retriever}: These methods use pre-defined rules or templates to extract paths from the graph. GenTKG~\cite{liao2024gentkg} considers a temporal KG, where they first extract logical rules from the KG, and use the top k rules to extract paths in a given time interval for the seed entities. Both \cite{dai2024counter, luo2023chatkbqa} generate queries using SPARQL, which are then used to retrieve import paths. KEQING~\cite{keqing} decomposes the original query into $k$ sub-queries using using an LLM fine-tuned via LoRA~\cite{hulora}. For each sub-query, they find the most similar question templates, which are predefined. For each template, they further pre-define a set of logical chains, which are then used to extract matching paths for the seed entities in the sub-query from the KG. 
    
    \item {\bf GNN-based retriever}: GNN-RAG~\cite{mavromatis2024gnn} trains a GNN for the retrieval task. A separate round of message passing is done for each query $q$, which is incorporated in the message computation along with the relation and entity representations. The GNN is then trained as in the node classification task, where the correct answer entity for $q$ has a label of 1 and 0 otherwise. During inference, the entities with probability above some threshold are treated as candidate answers, and the shortest path from the seed entity is extracted. \citet{gnn_llm} use a conditional GNN~\cite{conditional_mpnn} for retrieval, where for each query, only the seed entity (they assume there is only one) is initialized to a non-zero representation based on the LLM-encoded query. They then run $L$ rounds of message passing where after each layer $l$ only the top-K new edges are kept, resulting in a set of entities $C_q^l$. This is determined by a learnable attention weight, which prunes the other edges from the graph. The final set of candidate entities is the union of candidate entities at each layer $l$, $C_q = \cup_{l=1}^L C_q^l$. It is optimized in a similar manner to~\cite{mavromatis2024gnn}. For each candidate entity, they retrieve an evidence chain by backtracking from the entity until it reaches the seed entity, choosing those edges with the highest attention weight.
    REANO~\cite{fang2024reano} initializes the entity and relation representations via the mean-pooled representations of all mentions of that entity/relation in the texts, encoded by T5. They then run a GNN, which includes an attention weight that considers the relevance of a given triple to the original question (also encoded by T5). After running the GNN, they retrieve the top K triples in the KG that are most relevant to the question, where relevance is defined via the dot product between the triple encoded by the GNN and the question.  

    \item {\bf Similarity-based retriever}: STaRK~\cite{wu2024stark} considers the vector similarity of the query to each entity. Each entity embeds both the textual and relational information together. They further consider multi-vector similarity, where the entities are encoded using multiple vectors. This is done by chunking the textual and relational information of each entity, with each chunk being embedded into its own vector. Both REALM~\cite{zhu2024realm} and EMERGE~\cite{zhu2024emerge} extract the entities most similar to the query. While REALM only retrieves the entities themselves, EMERGE further retrieves the 1-hop subgraph around each entity.
    

    \item {\bf Relation-based retriever}: \citet{kim2023kg} propose a general framework for reasoning on KGs using LLMs. They first use an LLM to segment the original query into a set of $i \in I$ sub-sentences, where each sub-sentence $S_i$ has an associated set of entities $\mathcal{E}_i$. For each sub-sentence, they further use an LLM to retrieve the top-$k$ most relevant relations $\mathcal{R}_{i, k}$. Given the set of $k$ relations, for each sub-sentence, they retrieve all triples that contain a relation in $\mathcal{R}_{i, k}$ and whose entities are in $\bigcup_{i \in I} S_i$. GenTKGQA~\cite{gao2024two} focuses on temporal KG QA. Like~\cite{kim2023kg}, they retrieve the top-k relations for the query. They then retrieve all facts that contain one of the top k relations and satisfy the temporal constraints. 

    \item {\bf Fusion-based retriever}: These techniques consider a combination of different retrieval techniques. Mindmap~\cite{wen2023mindmap} considers extracting some paths $\leq k$ hops from the seed entity and the 1-hop subgraph of each seed. These two extracted components are combined into one subgraph. DALK~\cite{dalk}
    uses a procedure similar to Mindmap, where they extract both paths and the 1-hop subgraph around each seed entity. However, they argue that this procedure often results in the retrieval of redundant or unnecessary information. To remove these facts, they use an LLM to rank the retrieved facts given both the original question and the subgraph. Only the Top-k most relevant facts are kept.
    UniOQA~\cite{unioqa} considers two branches for retrieving. The first is a translator, which is a fine-tuned LLM that generates the answer in a CQL format. the second is a searcher that retrieves the 1-hop subgraph around the seed entities. When determining the answer, answers from the translator are prioritized over those from the searcher. KG-Rank~\cite{kgrank} considers ranking all triples in the 1-hop neighborhood of the seed entities via the similarity of the relation to the query, the similarity of each triple to the encoded output of $a = \text{LLM}(q)$, and an MMR ranking~\cite{carbonell1998use} that uses the similarity score. Only the top-ranked triples are kept. GrapeQA~\cite{taunk2023grapeqa} extends~\cite{yasunaga2021qa} by further including a set of ``extra nodes'', which are the common neighbors of the entities retrieved via a path-based retriever. They further introduce a clustering-based method for pruning entities that may be irrelevant to the query. SubgraphRAG~\cite{li2024simple} considers both GNN and textual information. For the GNN, they consider initializing the node representations using a one-hot encoding to differentiate between seed entities and others. A GNN is then run for L layers, resulting in the final representation $s_v$ for a node $v$. To retrieve the relevant triples, they consider first concatenating the final node representations for each triple $(h, r, t)$ such that $z_{\tau} = [s_h, s_t]$. The probability of choosing this triple is then given by  $p(h, r, t) = \text{MLP}([z_q, z_h, z_r, z_t, z_{\tau}])$, where $z_q$, $z_h$, $z_r$, $z_t$ are the encoded textual representations of the query and the triple $(h, r, t)$, respectively. Only the top K triples are chosen.
    
    \item {\bf Agent-based retriever}: These techniques use LLM agents to retrieve facts from the KG. KnowledGPT~\cite{wang2023knowledgpt} defines a set of tools for searching over a KG. Given a query, they generate a piece of code to search over the KG that considers the seed entities. The code is then executed over the KG to find the correct answer. KG-Agent~\cite{kgagent} focuses on fine-tuning an LLM to generate the SQL code for retrieving the correct answer. Using a set of tools, they extract a set of paths that contain the seed entities. KnowAgent~\cite{zhu2024knowagent} first identifies the relevant actions for the query via a planning module. Using these actions, they then generate a set of paths that are used for generation. 

    {\bf Other retrievers}: KICGPT~\cite{kicgpt} is concerned with the task of knowledge graph completion, where given a partial fact $(h, r, *)$, we want to predict the correct entity $\hat{e}$. KICGPT retrieves the entities by first scoring all possible entities using a traditional KG embedding score function. That is, for a score function $f(\cdot)$ and a partial fact $(h, r, *)$, they compute the set of scores $\{f(h, r, e) \; \forall e \in \mathcal{V} \}$. They use RotatE~\cite{sunrotate} for the function  $f(\cdot)$, a popular approach. Only the top $k$ entities by score are retrieved. To supplement their knowledge, they also retrieve all triples with (a) the same relation as the query and (b) all triples that contain the entity $h$ in the query. These are referred to as the analogous and supplement triple pools, respectively.

\end{itemize}


\subsection{Organizer}

In this subsection, we describe how the retrieved knowledge is organized for generation. More concretely, this is how the information is formatted when given to the generator. Note that not every method necessarily has an explicit organizer. We summarize the common methods below:
\begin{itemize}[leftmargin=*]
    \item {\bf Tuple-based organizer}: These methods consider each piece of retrieved information as an ordered triple. For example, it would include a triple in the generation prompt as ``$(\text{entity} \: 1, \text{relation} \: 1, \text{entity} \: 2)$''. Similarly, a path of length $m$ is given by ``$(\text{entity} \: 1, \text{relation} \: 1, \text{entity} \: 2, \text{relation} \: 2, \cdots, \text{entity} \: m)$''. The entities and relations are usually represented either as their names or IDs. Each triple or path is usually listed on a separate line. Many works append the retrieved paths to the original query as additional context \cite{sunthink, jiang2024hykge, mavromatis2024gnn, sanmartin2024kg, luoreasoning, choudhary2023complex, gnn_llm, wang2024reasoning, zhu2024knowagent}. Other works that retrieve facts instead of paths operate in a similar manner, where instead they append the triples~\cite{niu2024mitigating, zhu2024realm, dai2024counter, liao2024gentkg, kgrank, jiang2023structgpt, kim2023kg, gao2024two}. Some methods~\cite{zhu2024emerge, tian2024graph} consider only including the retrieved entities as the context. Given a set of facts, KG-R3~\cite{pahuja2023retrieve} first lists all entities and then relations, i.e., ``$(\text{entity} \: 1, \text{entity} \: 2, \cdots , \text{entity} \: m, \cdots, \text{relation} \: 1, \cdots, \text{relation} \: m-1)$''. 
    \citet{delile2024graph} consider a KG where each entity has an associated chunk of text. Each text chunk for an entity is considered as a different piece of information to be included in the context. Both \cite{tian2024graph, gao2024two} represent each entity and relation as an embedding, which is the combination of the LLM and GNN embedding. \citet{gnn_llm} further include the probability of each path containing the correct answer given by the GNN model. MVP-Tuning~\cite{mvp_tuning} considers combining multiple facts that share the same subject and relation to remove redundant information. That is, for a subject-relation pair $(\text{subject}, \text{relation})$, they denote the facts for $k$ possible objects as ``$\text{subject} \: \text{relation} \: \{\text{object} \: 1,  \cdots, \text{object} \: k\}$''. KG-Agent~\cite{kgagent} stores the current KG information and the historical reasoning programs in lists.

    \item {\bf Text organizer}: \citet{retrieve_rewrite_answer} verbalize the retrieved subgraph by passing each triple to the LLM and prompting it to convert it to a text representation. MindMap~\citet{wen2023mindmap} uses a similar procedure for subgraphs, where each is organized as a path before being passed to the LLM. Some methods use a set of pre-defined templates to verbalize the triples or paths \cite{knowledgenavigator, unioqa, kim2024causal}. \citet{keqing} experiment with verbalizing either via an LLM or pre-defined question templates, finding that LLM-based verbalizing works better for ChatGPT while template-based works better for LLaMA~\cite{touvron2023llama}. KICGPT~\cite{kicgpt} uses a combination of data preprocessing and LLM prompting to convert the triples to text. StaRK~\cite{wu2024stark} uses an LLM to synthesize each entity with its relational and textual information. Note that they use some pre-defined templates that depend on the specific task. CoTKR~\cite{wu2024cotkr} uses an LLM to summarize and then re-write a subgraph of facts for a question through a ``knowledge rewriter''. To train the rewriter, preference alignment is used, which optimizes the rewriter's output to match our preferred output. First, $k$ representations of the retrieved subgraph are produced, with ChatGPT choosing the best and work representations as the most and least preferred solutions.  
    
    \item {\bf Other organizer}: There are some exceptions to the previous classification. KnowledGPT~\cite{wang2023knowledgpt} represents the information in the form of a python class format. They also experiment with including additional information like the entity description and entity-aspect information.  

    \item {\bf Re-Ranking}: Some methods also {\it re-rank} the information in a specific order. This is done as the order of information can have a subtle impact on LLM performance. \citet{delile2024graph} order the text chunks of each entity based on the impact (measured by \# of citations of the parent paper) and the recency. \citet{dai2024counter} sort the triples by the relevancy score to the triple. \citet{choudhary2023complex} attempt to order the paths in a logical matter, such that for a given path, the subsequent paths build upon it. \citet{kgrank} re-rank the retrieved triples using a task-specific Cross-Encoder~\cite{jin2023medcpt}. STaRK~\cite{wu2024stark} considers re-ranking the retrieved entities using an LLM. The LLM is given the relational and textual information of each, and is asked to give it a score from 0 to 1, which is then used for re-ranking. GenTKG~\cite{liao2024gentkg} orders the paths by the time they occurred, further including the time with each. KICGPT~\cite{kicgpt} ranks all entities using the score of the KG embedding score function, keeping only the top $k$ entities. KICGPT re-rank the entities using in-context learning, where they prompt the LLM with examples from the analogy and supplement pool, as prior knowledge to aid the LLM in how to re-rank the entities. 
    
\end{itemize}

\subsection{Generator}

In this section we describe how the retrieved and organized data is used to generate a final response to the query. We categorize these generators according to the type of methods used to create these responses.

\begin{itemize}[leftmargin=*]
\item {\bf LLM-based generator}: The vast majority of works use a LLM to generate the response. The input to the LLM is the original query and retrieved and organized context, formatted using a specific template. The most commonly used LLMs include ChatGPT~\cite{chatgpt_paper}, Gemini~\cite{team2023gemini}, Mistral~\cite{jiang2023mistral}, Gemma~\cite{team2024gemma}, among others. For open-source models where the weights are publicly available, fine-tuning is sometimes used to modify the weights for a specific task \cite{unioqa, zhu2024knowagent}. This is often done through LoRA~\cite{hulora}, which allows for efficient fine-tuning.

\item {\bf GNN-based generator}: Some methods use graph neural networks (GNNs)~\cite{kipfgcn} to conduct the generation. \citet{yasunaga2021qa, taunk2023grapeqa, feng2020scalable} extract both the language and GNN embeddings for each potential answer (i.e., entity) conditional on the query. The probability of a single entity being the answer is then learnt based on the fusion of the two types of embeddings.



\item {\bf Other generators}: ~\citet{zhanggreaselm, yasunaga2022deep, hu2022empowering} formulate the prediction as a masked language modeling (MLM) problem. The goal is to predict the correct value (i.e., entity) for the masked token which answers the query. To do so, they fine-tune RoBERTa~\cite{liu2019roberta} language model. KG-R3~\cite{pahuja2023retrieve} scores the potential answer entities by performing cross-attention the representations of the query and each individual entity.
PullNet~\cite{pullnet} uses GraftNet~\cite{sun2018open} to score the different entities. \citet{gao2022graph} first selects the correct subgraph by computing the cosine similarity between the query and subgraph representations. For the subgraph with the highest similarity, it's fed to GraftNet~\cite{sun2018open} to select the most probable entity. REANO~\cite{fang2024reano} passes the encoded triples and their associated text passages to the T5 decoder. The task is framed as a classification problem, where the goal is to assign the highest probability to the triple with the correct answer.
\end{itemize}

\subsection{Resources and Tools} 

In this section, we list common tools and KGs that are used in graph RAG systems. For each, we give a brief description and a link to the project.

\subsubsection{Data Resources}
\begin{itemize}[leftmargin=*]
    \item {\bf Freebase}~\footnote{\href{https://developers.google.com/freebase}{https://developers.google.com/freebase}}~\cite{bollacker2008freebase} is an encyclopedic KG that contains a large variety of general and basic facts.
    \item {\bf ConceptNet}~\footnote{\href{https://conceptnet.io/}{https://conceptnet.io/}}~\cite{speer2017conceptnet} is a semantic graph, where the links in the graph are used to describe the meaning of different words or ideas.
    \item {\bf WikiData}~\footnote{\href{https://www.wikidata.org/wiki/Wikidata:Main\_Page}{https://www.wikidata.org/wiki/Wikidata:Main\_Page}}~\cite{vrandevcic2014wikidata} is a crowdsourced knowledge base that functions as a structured analog to the Wikipedia encyclopedia.
\end{itemize}

\subsubsection{ Tools}
\begin{itemize}[leftmargin=*]
    \item {\bf Graph RAG}~\footnote{\href{https://github.com/microsoft/graphrag}{https://github.com/microsoft/graphrag}}~\cite{edge2024local} is an official open-source implementation of the Graph RAG~\cite{edge2024local} framework. It can further be installed via the {\it graphrag} python package.
    \item {\bf LangChain}~\footnote{\href{https://python.langchain.com/v0.1/docs/use\_cases/graph/constructing/}{https://python.langchain.com/v0.1/docs/use\_cases/graph/constructing/}} is an open-source framework for using LLMs with various components and applications, including RAG, where using RAG on KGs is supportive. 
\end{itemize}

